%! The Secant, Cosecant, and Cotangent Functions — Unit 7, Lesson 7.4 — Group Activity (Answer Key)
\documentclass[10pt]{article}
\usepackage{precalculus-article}
\usepackage{precalculus-key}   % pulls in -boxes; do NOT also load -boxes
\newcommand{\TallMath}[1]{$\displaystyle #1\rule[-1.4em]{0pt}{3.2em}$}

\begin{document}

\pageheader{Unit 7, Lesson 7.4}{Group Activity --- Answer Key}
\namepartnerperiod

\vspace{0.10in}

%----------------------------------------------------------------------
% Context
%----------------------------------------------------------------------
\begin{scenariobox}[Context: Reciprocal Trigonometric Functions]{plumlight}
The secant, cosecant, and cotangent functions are built from the reciprocals
of cosine, sine, and tangent. In this activity your group will compute values
of these functions, explore their domains and asymptotes, and investigate
how the graph of a reciprocal function relates to its parent. Choose
\textbf{one tier} based on your readiness level.
\end{scenariobox}

\vspace{0.08in}

%----------------------------------------------------------------------
% Tier R
%----------------------------------------------------------------------
\begin{tcolorbox}[colback=white, colframe=black!40,
    title=\textbf{Tier R --- Reinforce},
    fonttitle=\bfseries, arc=2mm, left=3mm, right=3mm, top=2mm, bottom=2mm]

\textbf{Directions:} Use the given values of $\sin\theta$ and $\cos\theta$ to
compute $\sec\theta$, $\csc\theta$, and $\cot\theta$. Write ``und'' if the
function is undefined at that angle.

\vspace{4pt}
\renewcommand{\arraystretch}{1.7}
\begin{tabular}{|c|c|c|c|c|c|}
\hline
$\theta$ & $\sin\theta$ & $\cos\theta$ & $\sec\theta$ & $\csc\theta$ & $\cot\theta$ \\
\hline
$0$ & $0$ & $1$ & \ans{1} & \ans{und} & \ans{und} \\
\hline
$\dfrac{\pi}{6}$ & $\dfrac{1}{2}$ & $\dfrac{\sqrt{3}}{2}$ & \ans{$\dfrac{2\sqrt{3}}{3}$} & \ans{2} & \ans{$\sqrt{3}$} \\
\hline
$\dfrac{\pi}{4}$ & $\dfrac{\sqrt{2}}{2}$ & $\dfrac{\sqrt{2}}{2}$ & \ans{$\sqrt{2}$} & \ans{$\sqrt{2}$} & \ans{1} \\
\hline
$\dfrac{\pi}{3}$ & $\dfrac{\sqrt{3}}{2}$ & $\dfrac{1}{2}$ & \ans{2} & \ans{$\dfrac{2\sqrt{3}}{3}$} & \ans{$\dfrac{\sqrt{3}}{3}$} \\
\hline
$\dfrac{\pi}{2}$ & $1$ & $0$ & \ans{und} & \ans{1} & \ans{0} \\
\hline
$\pi$ & $0$ & $-1$ & \ans{$-1$} & \ans{und} & \ans{und} \\
\hline
$\dfrac{3\pi}{2}$ & $-1$ & $0$ & \ans{und} & \ans{$-1$} & \ans{0} \\
\hline
$2\pi$ & $0$ & $1$ & \ans{1} & \ans{und} & \ans{und} \\
\hline
\end{tabular}

\vspace{8pt}
\textbf{1.}\quad List all angles in the table where $\sec\theta$ is undefined:

\vspace{4pt}
$\theta = $ \ans{$\pi/2$ and $3\pi/2$}

\vspace{6pt}
\textbf{2.}\quad List all angles in the table where $\csc\theta$ is undefined:

\vspace{4pt}
$\theta = $ \ans{$0$, $\pi$, and $2\pi$}

\vspace{6pt}
\textbf{3.}\quad For the angles where $\csc\theta$ is defined, is $|\csc\theta| \geq 1$
always? Circle: \quad \ans{Yes}

\vspace{6pt}
\textbf{4.}\quad At $\theta = \pi/4$, is $\cot\theta$ positive or negative?
\ans{positive}
At $\theta = 3\pi/4$, is $\cot\theta$ positive or negative?
\ans{negative}

\end{tcolorbox}

\vspace{0.08in}

%----------------------------------------------------------------------
% Tier A
%----------------------------------------------------------------------
\begin{tcolorbox}[colback=white, colframe=black!40,
    title=\textbf{Tier A --- Apply},
    fonttitle=\bfseries, arc=2mm, left=3mm, right=3mm, top=2mm, bottom=2mm]

\textbf{Directions:} Answer each question without relying on a table. Show
your reasoning using definitions.

\vspace{6pt}
\textbf{1.}\quad State the domain of each function by listing the values excluded
from the domain.

\vspace{4pt}
(a)\quad $\sec\theta$: \ans{$\theta = \pi/2 + n\pi$ for integer $n$}

\vspace{4pt}
(b)\quad $\csc\theta$: \ans{$\theta = n\pi$ for integer $n$}

\vspace{4pt}
(c)\quad $\cot\theta$: \ans{$\theta = n\pi$ for integer $n$}

\vspace{8pt}
\textbf{2.}\quad List all vertical asymptotes of each function in the interval $[0, 2\pi]$.

\vspace{4pt}
$\sec\theta$: \ans{$\theta = \pi/2,\; 3\pi/2$} \quad
$\csc\theta$: \ans{$\theta = 0,\; \pi,\; 2\pi$} \quad
$\cot\theta$: \ans{$\theta = 0,\; \pi,\; 2\pi$}

\vspace{8pt}
\textbf{3.}\quad Evaluate each expression.

\vspace{4pt}
(a)\quad $\sec(2\pi/3) = \dfrac{1}{\cos(2\pi/3)} = \dfrac{1}{-1/2} = $ \ans{$-2$}

\vspace{6pt}
(b)\quad $\csc(5\pi/6) = \dfrac{1}{\sin(5\pi/6)} = \dfrac{1}{1/2} = $ \ans{2}

\vspace{6pt}
(c)\quad $\cot(5\pi/4) = \dfrac{-\sqrt{2}/2}{-\sqrt{2}/2} = $ \ans{1}

\vspace{8pt}
\textbf{4.}\quad \ans{No. Since $|\cos\theta| \leq 1$ for all $\theta$, we have $|\sec\theta| = |1/\cos\theta| \geq 1$. The value $0.8$ is between $-1$ and $1$, so it is impossible for $\sec\theta = 0.8$.}

\end{tcolorbox}

\vspace{0.08in}

%----------------------------------------------------------------------
% Tier E
%----------------------------------------------------------------------
\begin{tcolorbox}[colback=white, colframe=black!40,
    title=\textbf{Tier E --- Extend},
    fonttitle=\bfseries, arc=2mm, left=3mm, right=3mm, top=2mm, bottom=2mm]

\textbf{Directions:} Investigate the relationship between a function and its reciprocal.

\vspace{6pt}
\textbf{1. Arcs of Secant.}

\vspace{4pt}
(a)\quad $\theta = $ \ans{$0$, $\pi$, $2\pi$}

\vspace{4pt}
(b)\quad $\sec\theta = $ \ans{$1$, $-1$, $1$}

\vspace{4pt}
(c)\quad \ans{At a vertex, $|\cos\theta| = 1$ so $|\sec\theta| = 1$ as well. Since $|\sec\theta| \geq 1$ always, the arc cannot dip below $1$ (or above $-1$) --- it can only curve away from the $\theta$-axis at each vertex, never crossing it.}

\vspace{6pt}
\textbf{2. Algebraic Verification.}

\vspace{4pt}
\ans{$\cot\theta = \dfrac{1}{\tan\theta} = \dfrac{1}{\;\dfrac{\sin\theta}{\cos\theta}\;} = \dfrac{\cos\theta}{\sin\theta}$, which equals $\dfrac{\cos\theta}{\sin\theta}$ by definition of division of fractions. $\square$}

\vspace{6pt}
\textbf{3. Near an Asymptote.}

\vspace{4pt}
\ans{As $\theta \to 0^+$, $\sin\theta \to 0^+$ and $\cos\theta \to 1$, so $\cot\theta = \cos\theta/\sin\theta \to 1/0^+ = +\infty$. The table confirms: at $\theta = \pi/6$, $\cot = \sqrt{3} \approx 1.73$; at $\theta = \pi/12$ (not in table), $\cot \approx 3.73$. As $\theta$ decreases toward 0, cot increases without bound.}

\end{tcolorbox}

\end{document}
